\titledquestion{Master Theorem}

Assume that $n$ is a power of $2$, and every basic operation other than a
recursive call takes $\Theta(1)$ time.

Consider the following procedure:

\begin{cpp}
void Mystery(int n){
    if(n <= 1) return;

    for(int r = 0; r < 3; r++){
        Mystery(n / 2);
    }

    for(int i = 0; i < n; i++){
        for(int j = 0; j < floor(sqrt(n)); j++){
            // perform a constant-time operation
        }
    }
}
\end{cpp}

Give tight asymptotic bounds in $\Theta(\cdot)$ notation.

\begin{parts}

%---------------------------------------------------------
\part[3]
Complete the recurrence for the running time of \texttt{Mystery}:

\[
T(n)
=
\underbrace{\rule{1.2cm}{0.15mm}}_{a}
T\left(
\frac{n}{\underbrace{\rule{1.2cm}{0.15mm}}_{b}}
\right)
+
\Theta\left(
\underbrace{\rule{2cm}{0.15mm}}_{f(n)}
\right).
\]

\begin{solution}

$$
a=3
$$

$$
b=2
$$

$$
f(n) = n^\frac 3 2
$$

$$
T(n) = 3T(n/2) + \Theta(n^{3/2})
$$

\end{solution}


%---------------------------------------------------------
\part[4]
Apply the Master Theorem to the recurrence in part (a). State the applicable
case and give a tight asymptotic bound for $T(n)$.

\begin{solution}


$$
\log_b a = \log_2 3 \approx 1.585 > 1.5
$$

$$
T(n) = \Theta(n^{\log_2 3})
$$

\end{solution}


%---------------------------------------------------------
\part[4]
Use a recursion tree to analyze the recurrence in part (a).

\begin{subparts}

\subpart
At level $i$, find the number of subproblems, the size of each subproblem,
and the total number $W_i$ of constant-time operations performed by the
nested loops.

\subpart
Find the depth of the recursion tree and determine whether $W_i$ increases,
decreases, or remains asymptotically unchanged as $i$ increases.

\subpart
Derive a tight asymptotic bound for $T(n)$ from the recursion tree.

\end{subparts}

\begin{solution}

In layer $i$, there are $3^i$ subproblems, each of size $n/2^i$.

The operation number of layer $i$ is $W_i = 3^i \cdot c\left(\frac{n}{2^i}\right)^{3/2} = \Theta\left(n^{3/2} \left(\frac{3}{2\sqrt{2}}\right)^i\right)$.

The depth of recursion tree is $\log_2 n$.

$$
\frac{W_{i+1}}{W_i} = \frac{3}{2^{3/2}} = \frac{3}{\sqrt{8}} > 1
$$

$W_i$ increases as $i$ increases.

$$
\sum_{i=0}^{\log_2 n - 1} W_i = c \cdot n^{3/2} \sum_{i=0}^{\log_2 n - 1} \left(\frac{3}{2^{3/2}}\right)^i = c \cdot n^{3/2} \cdot \frac{\left(\frac{3}{2^{3/2}}\right)^{\log_2 n} - 1}{\frac{3}{2^{3/2}} - 1}
$$

$$
\left(\frac{3}{2^{3/2}}\right)^{\log_2 n} = \frac{3^{\log_2 n}}{\left(2^{\log_2 n}\right)^{3/2}} = \frac{n^{\log_2 3}}{n^{3/2}}
$$

$$
\sum_{i=0}^{\log_2 n - 1} W_i = \Theta\left(n^{3/2} \cdot \frac{n^{\log_2 3}}{n^{3/2}}\right) = \Theta\left(n^{\log_2 3}\right)
$$

$$
T(n) = \Theta\left(n^{\log_2 3}\right)
$$
\end{solution}


%---------------------------------------------------------
\part[4]
For $p>0$, analyze
\[
T_p(n)=4T_p(n/2)+\Theta(n^p).
\]

Give the Master Theorem case and the tight asymptotic bound for each possible
range of $p$.

\begin{solution}

$$
\log_b a = \log_2 4 = 2
$$

When $0 < p < 2$, $T_p(n) = \Theta(n^2)$.

When $p = 2$, $T_p(n) = \Theta(n^2 \log n)$.

When $p > 2$, $T_p(n) = \Theta(n^p)$.

\end{solution}


%---------------------------------------------------------
\part[2]
Consider the recurrence
\[
R(n)=R(n/2)+R(n/4)+\Theta(n^{3/2}).
\]

Can the Master Theorem be applied directly to this recurrence?
If not, use appropriate upper and lower bounds together with the Master
Theorem to determine a tight asymptotic bound for $R(n)$.

\begin{solution}

Master Theorem cannot be applied directly to this recurrence. Master Theorem requires that the the sub-problems have the same size, but in this recurrence, the sub-problems reduce to different sizes.

$$
R(n) = R(n/2) + R(n/4) + \Theta(n^{3/2}) \le 2R(n/2) + \Theta(n^{3/2})
$$

$$
T_{\text{upper}}(n) = 2T_{\text{upper}}(n/2) + \Theta(n^{3/2})
$$

$$
T_{\text{upper}}(n) = \Theta(n^{3/2})
$$

$$
R(n) = R(n/2) + R(n/4) + \Theta(n^{3/2}) \ge 2R(n/4) + \Theta(n^{3/2})
$$

$$
T_{\text{lower}}(n) = 2T_{\text{lower}}(n/4) + \Theta(n^{3/2})
$$

$$
T_{\text{lower}}(n) = \Theta(n^{3/2})
$$

Therefore, $R(n) = \Theta(n^{3/2})$.
\end{solution}

\end{parts}